\section{Preparing a consistent tuple for compression}
\label{sec:compression-local}

Our aim is to encode a long consistent tuple using almost $4\log q$ units
of information per flag.  Lemma~\ref{lem:stream-entropy} says that a tuple
selected from the stream requires an additional $\eta\log\log q$ per flag.
The present section prepares one compression step: two scans give small
joint domains, a short exposure makes representative flags nearly
independent of most targets, and the geometry of low-conflict pairs
reduces the problem to two endpoint supports.

Throughout this section and the next, fix $0<\eta<1/10$ and put
\[
 \beta=\eta/10^7,\qquad \sigma=\log q,\qquad
 k=\lfloor q\sigma^{1+\eta}\rfloor.
\]
All assertions are for sufficiently large primes $q$, with the threshold
allowed to depend on $\eta$.  A stage starts with a consistent tuple
$F=(F_i)_{i=1}^{\ell}$ of deterministic length $\ell=\Theta(k)$, in stream
order or in the globally reversed and swapped order.  Its underlying
stream marginal has density bounded by a constant relative to the iid
stream law, and the rectangle-occupancy event of
Lemma~\ref{lem:rectangles} holds almost surely.  Constants may depend on the stage, but the number of stages
will depend only on $\eta$.

A discrete context $\mathsf C$ specifies, for each tuple position, a
decoder-known set of flags containing its actual value.  The deterministic
parameters $\Lambda,\Delta\ge0$ satisfy
\begin{equation}\label{eq:stage-input}
 H(\mathsf C)\le\Lambda,\qquad
 |\text{each slot domain}|\le Cq^4e^\Delta.
\end{equation}
Set
\begin{equation}\label{eq:stage-parameters}
 D=\sigma^\beta(1+\Lambda/k+\Delta\sigma^{-\eta}),\qquad
 K=D\sigma^{3\beta},\qquad K_*=D\sigma^{6\beta}.
\end{equation}
For now assume
\begin{equation}\label{eq:stage-range}
 \sigma^\beta\le D\le\sigma^{1-\eta/2}.
\end{equation}
In particular $K_*=o(\sigma)$ and $K=o(K_*)$.  We use the parameters
$L,R,P$ of Theorem~\ref{thm:predictor}, so that
$L=D\sigma^{8\beta}$, $R$ is even and belongs to
$[\sigma^\beta,2\sigma^\beta+2]$, and $P=LR$.
The step will retain a fixed positive fraction of the tuple and give it
a new context with controlled entropy and slot-domain sizes.
Proposition~\ref{prop:compression-stage} states the quantitative output.

\subsection{Two marking scans}

Counting ordered independent sets by steps with few options or a large
reduction in the remaining domain goes back to Alon and R\"odl
\cite[Theorem~2.1 and Lemma~2.2]{AlonRodl2005}.
The following scan adapts the marking argument of
Brada\v c~\cite[Claim~2.13]{Bradac2026}, freezing its state between
expensive steps and recording its joint option counts for later entropy estimates.

\begin{lemma}[Marking]\label{lem:marking}
There is a message $\Omega$, consisting of constant-alphabet masks and
$O(q\sigma)$ full flags, with the following properties.
Given $(\mathsf C,\Omega)$, every unspecified flag belongs to a known set
of at most $C_1q^4$ flags, where $C_1$ is an absolute constant.  Moreover,
\begin{equation}\label{eq:marking-cost}
 H(\Omega\mid\mathsf C)
 \le O(k)+(4\sigma+\Delta+O(1))\EE h,
\end{equation}
where $h\le Cq\sigma$ is the number of specified flags.
Put $\theta_0=(\mathsf C,\Omega)$ and $J_0=\log(C_1q^4)$, and let
$F_{\rm cheap}$ be the tuple of unspecified flags. Its expected joint
entropy deficit satisfies
\begin{equation}\label{eq:deficit-budget}
 \EE\bigl[(\ell-h)J_0-H_{\theta_0}(F_{\rm cheap})\bigr]
 \le \mathcal B:=C(\Lambda+k+\Delta q\sigma).
\end{equation}
The same bound holds for any subset of unspecified positions determined
by $(\mathsf C,\Omega)$.
\end{lemma}

\begin{proof}
Scan forward, updating the state only at flags declared expensive.
For a point $y$ in the second point system, let
\[
 V(y)=\operatorname{span}\{b_i:a_i\perp y, i
       \text{ an earlier expensive position}\}.
\]
Let $U_j^0=\{y:\dim V(y)\le j\}$ and
$Z_j=\{y:\dim V(y)=j\}$, where dimensions are vector dimensions.
At an actual extension $(a,b)$, consistency implies $a\perp V(b)$.
Thus $r=\dim V(b)\le4$, and there are at most $2q^{4-r}$ choices for
$a$ once $b$ and $r$ are known.

Choose $l\le r$ maximizing $|Z_l|$, with a fixed rule for ties.  Declare
the flag popular-cheap if
\[
 |\{y\in Z_l:b\subseteq V(y)\}|\ge |Z_l|/(16q).
\]
Counting point memberships in the $l$-dimensional vector subspaces
$V(y)$ shows that there are at most $32q^l$ such $b$.  Hence their joint
flag count is at most $64q^4$.  If this condition fails, declare the flag
poor-cheap when
\[
 |\{y\in Z_l:a\perp y\}|\le |Z_l|/(8q).
\]
Lemma~\ref{lem:incidence} bounds the number of these $a$ by
$Cq^5/|Z_l|$.  The available $b$ belong to $U_r^0$, whose size is at
most $5|Z_l|$.  Applying incidence mixing to these two sets gives at most
$Cq^4$ flags: their product is $O(q^5)$, and both the main term and the
mixing error are $O(q^4)$.

Otherwise declare the flag expensive and perform the update.
At least $|Z_l|/(8q)$ points of $Z_l$ are hit by $a$, and fewer than
$|Z_l|/(16q)$ already have $b\subseteq V(y)$.  At least
$|Z_l|/(16q)\ge |U_l^0|/(80q)$ therefore leave $U_l^0$.
Each of the five nonempty sets $U_l^0$ can undergo only $O(q\sigma)$
such multiplicative reductions.  This bounds the number of updates.

Run the same scan after reversing the tuple and swapping the two point
systems.  Transmit rank/type masks for both scans and every flag
expensive in either scan.  These data reconstruct both states at every
position, and every remaining position passes both cheap scans.
The masks cost $O(k)$; the specified flags can be encoded using their
old context domains.  This proves~\eqref{eq:marking-cost}.

The unspecified positions and their number are known from $\theta_0$,
and the specified flags have been revealed. Thus
\[
 \EE H_{\theta_0}(F_{\rm cheap})
 =H(F\mid\mathsf C,\Omega)\ge H(F)-H(\mathsf C,\Omega).
\]
Using Lemma~\ref{lem:stream-entropy} and~\eqref{eq:marking-cost} gives
\[
 \EE\bigl[(\ell-h)J_0-H_{\theta_0}(F_{\rm cheap})\bigr]
 \le C(\Lambda+k+\Delta q\sigma).
\]
The $4\sigma\EE h$ cost cancels the expected cap of the specified positions.
The deficit is nonnegative because the remaining domains have size at
most $e^{J_0}$.  Finally, omitted coordinates contribute entropy at most
their cap sum, proving the assertion about subsets.
\end{proof}

The masks also give useful separate endpoint bounds.  At a slot that is
poor-cheap in at least one orientation, fix such an orientation by a
priority rule and put
$u=\max(\sigma,\min(4\sigma,\log|Z_l|))$.  Its endpoint supports have
sizes at most
\begin{equation}\label{eq:reciprocal-caps}
 A_{\max}=Ce^{5\sigma-u},\qquad B_{\max}=Ce^u.
\end{equation}
The clamping uses the trivial ambient bound $Q_4=O(q^4)$ at the two
endpoints.  If the slot is popular-cheap in both scans, with forward
and reversed ranks $r,r'$, its endpoint caps and partner caps are
as follows.  Here $|B\mid A|$ is the largest number of permitted second
endpoints for a fixed first endpoint, and $|A\mid B|$ is defined symmetrically:
\begin{equation}\label{eq:high-caps}
 |A|\le Cq^{r'},\quad |B|\le Cq^r,\qquad
 |B\mid A|\le2q^{4-r'},\quad |A\mid B|\le2q^{4-r}.
\end{equation}
In a high class we use $A_{\max}=Cq^{r'}$ and $B_{\max}=Cq^r$.
For $r+r'\le5$, these imply~\eqref{eq:reciprocal-caps} with $u=r\sigma$.
Rank zero cannot be popular.

We divide the remaining positions into a fixed number of classes.
Reciprocal classes specify their orientation and either an integer band
$|u-r\sigma|\le K_*$, with $r\in\{1,2,3,4\}$, or an open band
\[
 (r-1)\sigma+K_*<u<r\sigma-K_*.
\]
For open bands $r\in\{2,3,4\}$.
The remaining, high classes specify $(r,r')$ with $r+r'>5$; whenever a
rank is four, orient the class so that $r=4$.
By pigeonholing, an event determined by $\theta_0$ and of probability
bounded below has a fixed class containing at least $c k$ positions.
Condition on that event and retain the first $\lfloor c k\rfloor$
positions.  Conditional laws at a fixed $\theta_0$ have not changed,
and~\eqref{eq:deficit-budget} changes only by a constant factor.
The deterministic parameters $D,K,K_*$ remain those chosen at the
start of the stage.  Conditioning here multiplies the expected deficit
bound by at most the reciprocal of a fixed positive probability, which
is absorbed in its stage constant.  It does not redefine $D$ from the
entropy of the conditioned old context.  The next stage will instead
use the deterministic message-length bound for its newly constructed
context.
If the orientation is swapped, reverse the retained order as well.
All subsequent lengths are deterministic.

Split this class into equal windows, discarding the final remainder.
Let $w$ be the resulting deterministic number of complete windows, before
any subsequent deletions.
Take window length, rounded to a multiple of four, of order
\begin{equation}\label{eq:window-lengths}
 m\asymp
 \begin{cases}
 qD\sigma^{\eta/2},&\text{open bands},\\
 q\sigma^{1+\eta/2},&\text{integer bands}.
 \end{cases}
\end{equation}
A high class uses one window of length $\Theta(k)$.
The early and late quarters of each window are representative blocks;
the middle half consists of targets.  In every case $\Theta(k)$
positions remain.

\subsection{Choosing a round before exposing its values}

Conditioning to reduce correlations is a standard entropy method;
compare~\cite[Lemma~4.5]{RaghavendraTan2011}.  We need the following
version with one representative per block.

In the following estimates, conditional entropies and mutual information
are evaluated at the indicated history.  The outer expectation averages
over that history and all fresh representative indices.

\begin{lemma}[Pre-round decoupling]\label{lem:decoupling}
After exposing a deterministic number of rounds of older representatives,
one can draw one fresh uniform unused position from each representative
block, without exposing their values, with the following properties.
Let $\theta$ consist of $\theta_0$ and all older indices and values.
For an unexposed position $i$, write
\[
 \delta_i=J_0-H_\theta(F_i),\qquad
 I_{ij}=I_\theta(F_i;F_j).
\]
Then
\begin{equation}\label{eq:decoupling-bounds}
 \EE\sum_i\delta_i\le\mathcal B,\qquad
 \EE\sum_i\max_{j\in\mathcal J_i}I_{ij}
 \le\frac{C\mathcal B}{qD\sigma^{3000\beta}},
\end{equation}
where $\mathcal J_i$ consists of fresh representatives outside $i$'s
own representative block, and contains all representatives for a middle
target.  The expected sums over fresh representative targets themselves
are at most $C/m$ times the corresponding all-target sums.
The fresh indices do not change the conditional flag marginals.
\end{lemma}

\begin{proof}
Draw one new unused position per block for each of
$T=\lfloor qD\sigma^{3000\beta}\rfloor$ rounds, exposing values after
each round.  The parameter choices give $T=o(m)$, so all unused block
sizes remain comparable to $m$.
For any conditional law of coordinates and any set $E$ of coordinates,
the chain rule gives
\begin{equation}\label{eq:tc-drop}
 \TC(F)-\EE[\TC(F_{E^c}\mid F_E)]
 =\TC(F_E)+\sum_{i\notin E}I(F_i;F_E).
\end{equation}
Here the new indices are independent of the values given past history.
The initial expected total correlation is at most $\mathcal B$, because
it is bounded by the joint deficit in $J_0$-sized domains.

For an unselected target, its maximum mutual information with a new
representative is bounded by $I(F_i;F_E)$.  For a target in a
representative block, the maximum over the other blocks is independent
of the draw in its own block.  Thus the selected targets' expected
contribution is at most $C/m$ times the all-target contribution.
Absorbing this $o(1)$ fraction and summing~\eqref{eq:tc-drop} over rounds
proves that some deterministic pre-round satisfies the second bound
in~\eqref{eq:decoupling-bounds}.

Revealing and removing $r$ coordinates cannot increase the expected
deficit of those remaining: the removed coordinates have entropy at
most $rJ_0$.  This proves the first bound.  Uniformity gives the same
$C/m$ estimate for selected deficits.  Choose the indicated pre-round
and stop before revealing its fresh values.  Its fresh indices are
independent of the remaining values given $\theta$, proving the final
assertion.
\end{proof}

\subsection{Low-conflict pairs and collision rectangles}

All conditional probabilities in this subsection are given a fixed
pre-round history $\theta$.  For a flag $(a,b)$ at slot $i$, denote its
joint and endpoint marginal probabilities by $p_i(a,b)$, $p_i^A(a)$,
and $p_i^B(b)$.  Call a first endpoint $a$ good if
\[
 p_i^A(a)\ge e^{-K}/A_{\max}\quad\text{and}\quad
 \Pr\bigl(p_i(a,b)>q^{-4}e^K\mid a\bigr)\le .02;
\]
define good second endpoints symmetrically.  In a high class impose
also the upper cutoffs
\[
 p_i^A(a)\le e^Kq^{-r'},\qquad p_i^B(b)\le e^Kq^{-r}.
\]
The mass of bad endpoints at slot $i$ is
\begin{equation}\label{eq:bad-endpoint-mass}
 O\bigl(e^{-K}+(\delta_i+1)/K\bigr).
\end{equation}
Indeed, the negative part of $\log(p_i(F_i)q^4)$ has bounded expectation
by the support cap $C_1q^4$, so its positive part has expectation at most
$\delta_i+O(1)$.  Markov's inequality proves the light-atom condition.
The lower cutoff loses at most $e^{-K}$.  For the high upper cutoff,
the partner cap gives
$H_\theta(A_i)\ge r'\sigma-\delta_i-O(1)$, and the same negative-part
argument with the marginal cap $Cq^{r'}$ applies; similarly for $B_i$.

Fix $\rho=.005$.  For a first endpoint $a$ define a vector subspace in
the second system by
\[
 \mathcal K_i(a)=
 \bigcap_{\Pr(b\not\perp z\mid a)\le\rho} z^\perp.
\]
Define $\mathcal L_j(y)$ in the first system by the analogous construction
from the conditional first endpoint at second endpoint $y$.
A basis of the span of the tests has at most five members; hence
\begin{equation}\label{eq:core-capture}
 \Pr(b\in\mathcal K_i(a)\mid a)\ge1-5\rho,
 \qquad a\perp\mathcal K_i(a),
\end{equation}
with the analogous assertions for $\mathcal L_j(y)$.
At a good $a$, a captured light conditional atom is at most
$q^{-4}e^{2K}A_{\max}$, and these atoms have total mass at least $.955$.
Counting points in a vector subspace therefore gives
\begin{equation}\label{eq:core-dimension}
 \dim\mathcal K_i(a)\ge
 \left\lceil5-\frac{\log A_{\max}+3K}{\sigma}\right\rceil,
\end{equation}
and the corresponding bound using $B_{\max}$ for $\mathcal L_j(y)$.

\begin{lemma}[Low-conflict geometry]\label{lem:low-conflict}
For $i<j$, write the two flags as $(a,b),(c,y)$ and compare their actual
joint law with the independent product of their conditional marginals.
In a reciprocal class, or the high class $(3,3)$, there are nonnegative
numbers $c_{ij}$ such that
\begin{equation}\label{eq:low-conflict}
 \Pr_{\mathrm{prod}}(a,y\text{ good},\ a\perp y)
 \le C(I_{ij}+c_{ij}),
 \qquad \sum_{i<j}c_{ij}\le C\sigma k.
\end{equation}
In open reciprocal bands one can take $c_{ij}=0$.
\end{lemma}

\begin{proof}
The conflict event $a\perp y$, $c\not\perp b$ has probability zero
under the actual joint law.  Coarsening relative entropy to that event
shows that its product probability is at most $I_{ij}$.
Except for hit mass at most $I_{ij}/\rho^2$, its reverse-miss probability
conditional on $(a,y)$ is at most $\rho^2$.
Markov's inequality then shows that at least $1-\rho$ of $c\mid y$
are tests defining $\mathcal K_i(a)$.  Every point of $\mathcal K_i(a)$
is consequently a test defining $\mathcal L_j(y)$, so
\begin{equation}\label{eq:core-containment}
 \mathcal L_j(y)\subseteq\mathcal K_i(a)^\perp.
\end{equation}
On such endpoints, both memberships
$b\in\mathcal K_i(a)$ and $c\in\mathcal L_j(y)$ hold under the product
law with probability at least $1-10\rho$.

In an open band,~\eqref{eq:core-dimension} gives dimensions at least
$r$ and $6-r$, contradicting~\eqref{eq:core-containment}.
In an integer band the bounds are $r$ and $5-r$, so equality holds in
the orthogonal complement.  Let $E_{ij}$ require good endpoint incidence,
the containment, and both core memberships.  Both flags then belong to
$\mathcal R(\mathcal K_i(a))$.

For $(3,3)$, both core dimensions are at least two.  Put
$V=\mathcal L_j(y)^\perp$.  Then
$\mathcal K_i(a)\subseteq V$ with dimension difference at most one.
If $y\in\mathcal K_i(a)$, both flags belong to
$\mathcal R(\mathcal K_i(a))$.  Otherwise
$V=\operatorname{span}(\mathcal K_i(a),y)$, and $a\perp y$ shows that
both flags belong to $\mathcal R(V)$.
The first rectangle is determined by the earlier slot, the second by
the later slot.  Lemma~\ref{lem:rectangles} thus bounds the number of
charged partners of any slot by $C\sigma$, pointwise for every actual
tuple.  Set $c_{ij}=\Pr_{\mathrm{joint}}(E_{ij}\mid\theta)$ and take
conditional expectations to obtain its sum bound.

For any pair event $E$, the entropy variational inequality with test
$-1_E$ gives
\[
 I_{ij}\ge-\Pr_{\mathrm{joint}}(E)
   -\log\bigl(1-(1-e^{-1})\Pr_{\mathrm{prod}}(E)\bigr).
\]
Therefore $\Pr_{\mathrm{prod}}(E)\le
C(I_{ij}+\Pr_{\mathrm{joint}}(E))$.  Apply this to $E_{ij}$ and use
the product core-capture bound to prove~\eqref{eq:low-conflict}.
\end{proof}

Extend $c_{ij}$ symmetrically, setting it to zero in the other high
classes.  For a fresh target $i$, put
\[
 g_i=\max_{j\in\mathcal J_i}(I_{ij}+c_{ij}).
\]
Call the index good if
\begin{equation}\label{eq:good-index}
 \delta_i\le D\sigma^\beta,\qquad
 g_i\le\sigma^{-2000\beta}/q.
\end{equation}
The expected fractions of bad indices and bad representatives are
$o(1)$.  For mutual information use Lemma~\ref{lem:decoupling} and
$\mathcal B/k=O(D\sigma^{-\beta})$.  For collisions, uniform block
sampling and~\eqref{eq:low-conflict} give expected total at most
$C\sigma k/m$; collision windows have
$m\ge cq\sigma^{1+\eta/2}$, so Markov loses at most
$C\sigma^{2000\beta-\eta/2}=o(1)$.  Selection in a target's own block
is independent of its other-block sum, giving the representative estimate
as before.  At every good index,~\eqref{eq:bad-endpoint-mass} is $o(1)$.

\subsection{Eliminating high classes}

\begin{lemma}[High ranks]\label{lem:high-ranks}
The retained class cannot be high.  In the rank-four case, the stated
assumptions instead directly produce an encoding contradicting
Lemma~\ref{lem:stream-entropy}.
\end{lemma}

\begin{proof}
First suppose the class is $(3,3)$.  Choose a good early representative
and a good middle target, which exist for some histories and fresh draws
by the preceding estimates.  Restrict the representative's first
marginal to good endpoints.  Its total mass is $1-o(1)$ and its atoms
are at most $e^K/q^3$.  Weighted incidence variance shows that all but
$Cq^2e^K$ second endpoints have incident mass at least $.3/q$.
The exceptional set has mass at most $Cq^{-1}e^{2K}=o(1)$ in the good
target's second marginal.  This contradicts~\eqref{eq:low-conflict}
and~\eqref{eq:good-index}.

Now orient so that $r=4$ and $r'\ge2$.  Draw
$h=\lceil q\sigma^\beta\rceil$ independent full flags from a good
early representative's marginal conditioned on both endpoints being
good.  These draws are independent of the actual targets given history.
Weighted variance shows that all but $o(1)$ of a good middle target's
second-endpoint mass has sample first-endpoint hit probability at least
$.3/q$: the exceptional set has size $Cq^{5-r'}e^K$, while its good
atoms are at most $e^K/q^4$.

For such a target endpoint $y$, start with its span and add the second
endpoints of hit samples.  While the span has vector dimension below
four, it contains at most $Cq^2$ points, whose sample mass is at most
$Cq^{-2}e^K=o(1/q)$.  Each trial therefore has conditional probability
at least $.2/q$ of increasing the dimension.  Binomial domination shows
that the span reaches dimension four with failure probability $o(1)$.
The probability of any conflict with the actual target is at most
$ChI_{ij}=o(1)$; conditioning the representative on good endpoints
changes its density only by a bounded factor.

In expectation only $o(k)$ middle targets are lost, also counting a bad
representative and bad target indices as losses.  With probability
bounded below retain a deterministic $\Theta(k)$ covered middle targets
in their original order.  Transmit the sample list, at cost
$O(h\sigma)=o(k)$.  For each possible $y$, the decoder forms the span
of $y$ and the second endpoints of its hit samples.  If its dimension
is at least four, at most one first endpoint annihilates it.  Every
covered actual target belongs to this common domain, of size at most
$Q_4$.  No old context is needed.  The resulting tuple has entropy at
most $4\sigma\ell+O(k)$, contradicting Lemma~\ref{lem:stream-entropy}
after the constant-probability conditioning.
\end{proof}

\subsection{Nearly uniform endpoint supports}

Only reciprocal classes remain.  Their endpoint caps are
\eqref{eq:reciprocal-caps}, with a history-known number $u_i$ at each
slot.

\begin{lemma}[Endpoint levels]\label{lem:flat-levels}
At every good index there are independent auxiliary choices of first
and second endpoint supports, each consisting of good endpoints, such
that a chosen support $A$ satisfies
\[
 \tfrac12 A_{\max}e^{-K}\le |A|\le A_{\max},
 \qquad \EE\frac{1_{\{a\in A\}}}{|A|}\le C p_i^A(a),
\]
and similarly for $B$.  For any two ordered good representatives in
distinct blocks,
\begin{equation}\label{eq:u-monotonicity}
 u_j\le u_i+CK\qquad(i<j).
\end{equation}
\end{lemma}

\begin{proof}
Partition each endpoint marginal into levels of
$\lfloor-\log p\rfloor$.  The entropy of this integer-valued level is
$O(\log\sigma)$: its mean is $O(\sigma)$, and comparison with a geometric
distribution of that mean gives entropy at most
$\log(e(1+O(\sigma)))$.
For a first level set $A$, put
$\operatorname{loss}_A=\log(A_{\max}/|A|)$, and define the second loss
similarly.  Conditional on both levels, incidence mixing gives
\[
 \log|\{(a,b)\in A\times B:a\perp b\}|
 \le4\sigma+O(1)
       -\tfrac12(\operatorname{loss}_A+\operatorname{loss}_B).
\]
The chain rule consequently bounds the expected total loss by
$O(\delta_i+\log\sigma+1)$.  A level is permissible when its loss is
at most $K$ and at least half its members are good.  Since probabilities
within a level differ by at most $e$, the total probability of permissible
levels is $1-o(1)$ by~\eqref{eq:bad-endpoint-mass} and Markov.
Choose a permissible level by its marginal probability, normalized over
permissible levels, and take its good subset.  This proves both the size
and averaged pointwise bounds.

For $i<j$, independent choices of a first support at $i$ and second
support at $j$ have expected incidence density at most
$C\sigma^{-2000\beta}/q$ by~\eqref{eq:low-conflict} and averaged
pointwise domination.  If $u_j-u_i>CK$ with a sufficiently large $C$,
every such pair has product at least
$c q^5\exp(u_j-u_i-2K)$; mixing forces density at least $c/q$.
This contradiction proves~\eqref{eq:u-monotonicity}.
\end{proof}

Discard windows with a bad representative.  In an open band discard
also windows with $u_{\rm early}-u_{\rm late}>K_*$.  Along the remaining
ordered representative list before this last deletion, the sum of
downward changes is at most $C\sigma+CKw$, by~\eqref{eq:u-monotonicity}.
Thus these extra deletions cost at most
$C(\sigma+Kw)/K_*=o(w)$ windows, since
$w\asymp\sigma^{1+\eta/2}/D$.  In an integer band the variation is
already at most $2K_*$.

At every surviving window independently choose a permissible first
support $A$ from its early representative and a permissible second
support $B$ from its late representative.  Then
\begin{equation}\label{eq:pivot-support-input}
 |A||B|\ge q^5e^{-CK_*}.
\end{equation}
For the same window and every ordered cross pair, expected incidence
density is at most $C\sigma^{-2000\beta}/q$.  The same bound holds
against the corresponding good endpoint mass of every good middle
target on the appropriate temporal side.  These are averaged statements
over the independent support choices; no pointwise sparsity of every
chosen pair is being assumed.  They are precisely the inputs for the
tree construction in the next section.
